\documentclass[CJK]{beamer}
\usepackage{CJKutf8}
\usepackage{beamerthemesplit}
\usetheme{Malmoe}
\useoutertheme[footline=authortitle]{miniframes}
\usepackage{amsmath}
\usepackage{amssymb}
\usepackage{graphicx}
\usepackage{color}
\usepackage{slashed}
\usepackage{simplewick}
\graphicspath{{../figures/}}
\def\be{\begin{equation}}
\def\ee{\nonumber\end{equation}}
\def\bea{\begin{eqnarray}}
\def\eea{\nonumber\end{eqnarray}}
\def\ii{{\dot{\imath}}}
\def\bch{\begin{CJK}{UTF8}{gbsn}}
\def\ech{\end{CJK}}
\def\bex{\begin{minipage}{0.3\textwidth}\includegraphics[width=1in]{jugelizi.png}\end{minipage}\begin{minipage}{0.6\textwidth}}
\def\eex{\end{minipage}}
\def\chtitle#1{\frametitle{\bch#1\ech}}
\def\skipline{{\vskip0.1in}}
\def\skiplines{{\vskip0.2in}}
\def\lagr{{\mathcal{L}}}
\def\hamil{{\mathcal{H}}}
\def\vecv{{\mathbf{v}}}
\def\vecx{{\mathbf{x}}}
\def\veck{{\mathbf{k}}}
\def\vecp{{\mathbf{p}}}
\def\vecn{{\mathbf{n}}}
\def\vecA{{\mathbf{A}}}
\def\vecP{{\mathbf{P}}}
\def\vecsigma{{\mathbf{\sigma}}}
\def\hatJn{{\hat{J_\vecn}}}
\def\hatJx{{\hat{J_x}}}
\def\hatJy{{\hat{J_y}}}
\def\hatJz{{\hat{J_z}}}
\def\hatj#1{\hat{J_{#1}}}
\def\hatphi{{\hat{\phi}}}
\def\hatq{{\hat{q}}}
\def\hatpi{{\hat{\pi}}}
\def\vel{\upsilon}
\def\Dint{{\mathcal{D}}}
\def\adag{{\hat{a}^\dagger}}
\def\bdag{{\hat{b}^\dagger}}
\def\cdag{{\hat{c}^\dagger}}
\def\ddag{{\hat{d}^\dagger}}
\def\hata{{\hat{a}}}
\def\hatb{{\hat{b}}}
\def\hatc{{\hat{c}}}
\def\hatd{{\hat{d}}}
\def\hatN{{\hat{N}}}
\def\hatH{{\hat{H}}}
\def\hatp{{\hat{p}}}
\def\Fup{{F^{\mu\nu}}}
\def\Fdown{{F_{\mu\nu}}}
\def\newl{\nonumber \\}
\def\SIkm{\mathrm{km}}
\def\SIyr{\mathrm{yr}}
\def\SIGyr{\mathrm{Gyr}}
\def\SIeV{\mathrm{eV}}
\def\SIGeV{\mathrm{GeV}}
\def\SIm{\mathrm{m}}
\def\SIcm{\mathrm{cm}}
\def\SIJ{\mathrm{J}}
\def\SIs{\mathrm{s}}
\def\SIkg{\mathrm{kg}}
\def\SIg{\mathrm{g}}
\def\vece{\mathrm{e}}
\def\bmat#1{\left(\begin{array}{#1}}
\def\emat{\end{array}\right)}
\def\bcase#1{\left\{\begin{array}{#1}}
\def\ecase{\end{array}\right.}
\def\calM{{\mathcal{M}}}
\def\calT{{\mathcal{T}}}
\def\calR{{\mathcal{R}}}
\def\barpsi{\bar{\psi}}
\def\baru{\bar{u}}
\def\barv{\bar{\upsilon}}
\def\bmini#1{\begin{minipage}{#1\textwidth}}
\def\emini{\end{minipage}}
\def\qeq{\stackrel{?}{=}}
\def\torder#1{\mathcal{T}\left(#1\right)}
\def\rorder#1{\mathcal{R}\left(#1\right)}


\title{Quantum Field Theory I \\ Lesson 14 - Cross section}
\author{}
\date{}


\begin{document}

\begin{frame}
 
\begin{center}
\begin{Large}
\bch
量子场论 I 

{\vskip 0.3in}

第十四课 散射截面

\ech
\end{Large}
\end{center}

\vskip 0.2in

\bch
课件下载
\ech
https://github.com/zqhuang/SYSU\_QFTI

\end{frame}

\begin{frame}
\chtitle{散射截面是粒子的有效面积}
\bch
对散射问题，如果一开始就是取一个粒子为参照系，即$\vecp_1 = 0$。不妨取粒子2的运动方向为$z$轴，其动量为
$$\vecp_2 = \left(0, 0, \frac{m\upsilon}{\sqrt{1-\upsilon^2}}\right)$$
其中$\upsilon$是粒子2在这个参考系里的运动速度。
\skipline

散射截面$\sigma$是粒子2的“有效面积”，粒子2在总时间$T$内扫过的体积为$\sigma \upsilon T$。在总体积为$V$的情况下，粒子2“碰撞”到粒子1的概率为
$$\frac{\sigma \upsilon T}{V}$$


\ech
\end{frame}


\begin{frame}
\chtitle{和散射振幅联系起来}
\bch
另一方面，我们计算了发生散射的概率为
{\small
$$ \sum_{\vecp_3, \vecp_4} \left\vert\calM \frac{T}{V}  \delta(p_1+p_2-p_3-p_4) d^4k\right\vert^2 = \frac{1}{2!}\sum_{\vecp_3, \vecp_4} |\calM|^2 \frac{T^2}{V^2}  \delta(p_1+p_2-p_3-p_4) d^4k $$
}
因为是总概率，所以要对所有$\vecp_3$, $\vecp_4$求和。但是，“产生一个动量为$\vecp_3$的粒子和另一个动量为$\vecp_4$的粒子”和“产生一个动量为$\vecp_4$的粒子和另一个动量为$\vecp_3$的粒子”描述的是同一件事情。所以有了外面的$\frac{1}{2!}$因子。

要求能量动量守恒的因子$\delta(p_1+p_2-p_3-p_4) d^4k$，因为它满足能量动量守恒时为1，否则为零。所以它的平方总是等于自己。

\skipline 
令两种方法算出来的概率相等，即有
$$\sigma  = \frac{1}{2!}\sum_{\vecp_3, \vecp_4} |\calM|^2 \frac{T}{V\upsilon}  \delta(p_1+p_2-p_3-p_4) d^4k $$
\ech
\end{frame}


\begin{frame}
\chtitle{化为标准形式}
\bch
再利用$d^4k = \frac{(2\pi)^4}{VT}$，以及$d^3\vecp_3 = d^3\vecp_4 = \frac{(2\pi)^3}{V}$，就得到散射截面公式的标准形式
$$\sigma  = \frac{1}{2!(2\pi)^2} \int d^3\vecp_3 \int d^3\vecp_4\, \frac{|\calM|^2}{\upsilon}  \delta(p_1+p_2-p_3-p_4)  $$
注意式中的$\delta$函数是四维的。并且，如果产生的两个粒子是两种不同的粒子，则没有$\frac{1}{2!}$因子。

\ech
\end{frame}


\begin{frame}
\bch
问题是，只会小学数学的我们，真的能算出这个积分吗？
\ech
\end{frame}


\begin{frame}
\chtitle{先算最简单的}
\bch
先来看最简单的情况，假设只有四次耦合项$\frac{\lambda}{4!}\phi^4$。那么
$$|\calM|^2 = \frac{\lambda^2}{16\omega_1\omega_2\omega_3\omega_4}$$

代入散射截面公式，
{\small
\bea
\sigma &=& \frac{\lambda^2}{128\pi^2\upsilon \omega_1\omega_2} \newl
&& \times \int d^3\vecp_3 \int d^3\vecp_4\, \frac{1}{\omega_3\omega_4}\delta\left(\omega_1+\omega_2-\omega_3-\omega_4\right)\delta\left(\vecp_2 - \vecp_3 -\vecp_4\right) 
\eea
这里
$$\omega_1 = m,\  \omega_2 = \frac{m}{\sqrt{1-\upsilon^2}},\ \omega_3 = \sqrt{m^2+\vecp_3^2}, \ \omega_4 = \sqrt{m^2+\vecp_4^2}$$
}
 
\ech
\end{frame}


\begin{frame}
\chtitle{把$\vecp_4$积掉}
\bch
{\small 
可以先对$\vecp_4$积分，得到
\be
\sigma = \frac{\lambda^2}{128\pi^2 \upsilon\omega_1\omega_2} \int d^3\vecp_3 \, \frac{1}{\omega_3\omega_4}\delta\left(\omega_1+\omega_2-\omega_3-\omega_4\right)
\ee
当然，上式中$\omega_4$的含义发生了变化：
$$\omega_4 = \sqrt{m^2+(\vecp_2-\vecp_3)^2}$$
事实上，在取定的参考系里，$\omega_4\ge m$，又要能量守恒$\omega_1+\omega_2  =\omega_3+\omega_4$，所以$\omega_3\le \omega_2$。那么上面的积分仅在
$$|\vecp_3| \le |\vecp_2| $$
时有非零贡献。
}
\ech
\end{frame}


\begin{frame}
\chtitle{取球坐标系}
\bch
{\small
设$\vecp_3$的球座标系坐标为$(|\vecp_3|,\theta,\varphi)$, 则
$$\int d^3\vecp_3 = \int_0^{|\vecp_2|} |\vecp_3|^2 d|\vecp_3| \int_{-1}^1 d(\cos\theta) \int_0^{2\pi}d\varphi $$
因为$\sigma$的积分函数与$\varphi$无关，可以先对$\varphi$积分得到$2\pi$。
于是
$$\sigma = \frac{\lambda^2}{64\pi \upsilon\omega_1\omega_2}  \int_0^{|\vecp_2|} |\vecp_3|^2 d|\vecp_3| \int_{-1}^1 d(\cos\theta) \frac{1}{\omega_3\omega_4}\delta(\omega_1+\omega_2-\omega_3-\omega_4)  $$
这时$\omega_3$, $\omega_4$的含义为
$$\omega_3 = \sqrt{m^2+|\vecp_3|^2},\ \omega_4 = \sqrt{m^2+ |\vecp_2|^2 + |\vecp_3|^2 - 2|\vecp_2||\vecp_3|\cos\theta} $$ 
} 
\ech
\end{frame}

\begin{frame}
\chtitle{积掉$\cos\theta$}
\bch
{\small
先对$\cos\theta$进行积分，为此我们要计算出
$$\left\vert\frac{d(\omega_1+\omega_2-\omega_3-\omega_4)}{d (\cos\theta)}\right\vert = \frac{|\vecp_2||\vecp_3|}{\omega_4}$$ 
再利用积分公式
$$\int \delta\left(f(x)\right) dx = \sum_{x*:\ f(x*)=0} \frac{1}{|f'(x*)|}$$
就得到
$$\sigma = \frac{\lambda^2}{64\pi \upsilon\omega_1\omega_2|\vecp_2|}  \int_0^{|\vecp_2|}  d|\vecp_3| \frac{|\vecp_3|}{\omega_3} $$
注意：$\theta$有解的充分必要条件是$|\vecp_3|\le|\vecp_2|$（也就是$m\le\omega_3\le\omega_2$)，把$\theta$积掉后$|\vecp_3|$的积分上下限就必须确定了。
} 
\ech
\end{frame}

\begin{frame}
\chtitle{最后积掉$|\vecp_3|$}
\bch
最后利用$|\vecp_3|d|\vecp_3| =  \omega_3 d\omega_3$，就得到
\bea
\sigma &=& \frac{\lambda^2}{64\pi \upsilon\omega_1\omega_2|\vecp_2|}  \int_m^{\omega_2}  d\omega_3  \newl
&=& \frac{\lambda^2(\omega_2 - m)}{64\pi \upsilon\omega_1\omega_2|\vecp_2|} \newl
&=& \frac{\lambda^2}{64\pi  m^2\left(1+\frac{1}{\sqrt{1-\upsilon^2}}\right)} 
\eea
非相对论极限$\upsilon\ll 1$下，
$$\sigma \approx \frac{\lambda^2}{128\pi m^2}$$
\ech
\end{frame}

\begin{frame}
\chtitle{取质心参照系计算更方便}
\bch
{\small 
实际上，散射截面不一定要在一个粒子的静止参照系里计算。我们可以取质心参照系，设粒子1的速度为沿$z$轴向上$u$，粒子2的速度相反。在这个参照系里，若要能量动量都守恒则必须有
$$\vecp_3+\vecp_4 = 0$$
且
$$|\vecp_1| = |\vecp_2| = |\vecp_3| = |\vecp_4|$$
因为粒子之间相对速度为$2u$，上面开始几步步骤中需要把$\upsilon$替换为$2u$。直到我们在球坐标系里写下：
$$\sigma = \frac{\lambda^2}{64\pi (2u)\omega_1\omega_2}  \int_0^{\infty} |\vecp_3|^2 d|\vecp_3| \int_{-1}^1 d(\cos\theta) \frac{1}{\omega_3\omega_4}\delta(\omega_1+\omega_2-\omega_3-\omega_4)  $$
注意现在$\omega_3 = \omega_4 = \sqrt{|\vecp_3|^2+m^2}$，所以积分函数就不依赖于$\theta$，直接可以对$\theta$积分得到
$$\sigma = \frac{\lambda^2}{32\pi (2u)\omega_1\omega_2}  \int_0^{\infty} |\vecp_3|^2 d|\vecp_3|  \frac{1}{\omega_3\omega_4}\delta(\omega_1+\omega_2-\omega_3-\omega_4)  $$
}
\ech
\end{frame}

\begin{frame}
\chtitle{取质心参照系计算更方便}
\bch
{\small 
然后我们计算出
$$\left\vert\frac{d(\omega_1+\omega_2-\omega_3-\omega_4)}{d|\vecp_3|} \right\vert = \frac{2|\vecp_3|}{\omega_4}=\frac{2|\vecp_1|}{\omega_1}$$
利用$\delta$函数积分公式即得
$$\sigma = \frac{\lambda^2|\vecp_1|  }{128 \pi u \omega_1^3} $$
再利用$u\omega_1 =|\vecp_1|$，上式就化简为
$$\sigma = \frac{\lambda^2}{128 \pi \omega_1^2} = \frac{\lambda^2}{32 \pi E_{\rm tot}^2}  $$
其中$E_{\rm tot}=2\omega_1 $是质心系总能量。
}

\ech
\end{frame}


\begin{frame}
\chtitle{课堂讨论}
\bch
在两个参照系里算出来的结果一样吗？
\ech
\end{frame}

\end{document}


